\documentclass[a4paper,12pt]{ctexart}
\usepackage{xcolor}
\usepackage[top=1.8cm, bottom=1.8cm, left=1.8cm, right=1.8cm]{geometry}
\usepackage{array}
\linespread{1.35}

\newcommand{\ruby}[2]{#1\textsuperscript{#2}}

\newlength{\colA}\setlength{\colA}{0.06\textwidth}
\newlength{\colB}\setlength{\colB}{0.64\textwidth}
\newlength{\colC}\setlength{\colC}{0.26\textwidth}

\newcommand{\sectab}[3]{%
  \noindent
  \begin{minipage}[c]{\dimexpr\colA+\colB+2\tabcolsep\relax}
  \begin{tabular}{@{}p{\colA}p{\colB}@{}}
  \sideyuan  & {\large\bfseries #1} \\
  \sideyiwen & {\large #2} \\
  \end{tabular}
  \end{minipage}
  \hfill
  \begin{minipage}[c]{\colC}
  {\labelann\vspace{0.1cm}}#3
  \end{minipage}
  \par
}

\newcommand{\sideyuan}{\colorbox{blue!70!black}{\textcolor{white}{\normalsize\textbf{原文}}}}
\newcommand{\sideyiwen}{\colorbox{green!60!black}{\textcolor{white}{\normalsize\textbf{译文}}}}
\newcommand{\labelann}{\textcolor{orange!80!black}{\small\textbf{—— 注释 ——}}}
\newcommand{\refmark}[1]{\textsuperscript{\textcolor{red}{#1}}}

\title{\textcolor{blue!70!black}{\Huge\textbf{《左传》名篇选读}}}
\author{}
\date{}

\begin{document}
\maketitle
\thispagestyle{empty}

\section*{\textcolor{blue!70!black}{知人知世：关于《左传》}}

同学们，你们知道吗？在很久很久以前的春秋时期，距离今天已经有2500多年了。那时候有一位叫\textbf{左丘明}的历史学家，把他那个时代发生的大事都记录了下来，写成了一本书，这本书就叫《左传》。

《左传》的全名是《春秋左氏传》。它可不是一般的书哦！它既是一部厚重的历史书，也是一部精彩的文学书。书里的故事写得特别生动，有兄弟相争、有智者退敌、有以弱胜强……读起来就像看故事片一样过瘾。

那么问题来了——为什么我们读的《左传》原文和我们平时说的话不太一样呢？这是因为古人写文章用的是\textbf{文言文}。文言文是古代的书面语言，和当时人们日常说话也不太一样。

为什么文言文会这么简短呢？这里面有一个很有趣的原因：古时候没有纸，写字要用竹简——就是把竹子削成一片一片的，然后用绳子串起来。竹简又重又占地方，据说秦始皇每天批阅的奏章就有一百多斤重！所以古人写文章的时候，就尽量用最少的字表达最多的意思，慢慢就形成了文言文\textbf{简练}的特点。比方说，我们平时说"我的母亲"，在文言文里常常只用一个"母"字；我们说"你有没有吃饭"，古人可能就说"食否"两个字。这就是把口语中的多音节词简省成了单音节词。

另外，文言文里还有很多\textbf{虚词}，比如"之""乎""者""也""而""以""于"等等。这些词虽然没有实际的意义，但它们在句子中起到连接、停顿、强调等作用，是读懂文言文的关键。所以古人说"之乎者也"，指的就是这些虚词。掌握了这些虚词的用法，就像拿到了打开文言文大门的钥匙！

那么，我们怎么才能读懂《左传》这样的文言文呢？告诉大家几个小窍门。

第一，要学会\textbf{了解背景}。在读一篇文章之前，先了解一下作者是谁、当时发生了什么事，这样文章就容易弄懂了。这就叫"知人知世"。

第二，要\textbf{借助注释和译文}。我们这本书里，难懂的字词旁边都有注释——就像一本小词典。原文下面呢，还有现代文的翻译。你可以先读原文，再看注释，最后对照译文，慢慢就会发现古文也没那么难懂。

第三，要\textbf{多读多背}。文言文读多了，自然就会找到感觉，这就叫"语感"。有些优美的句子背下来，以后写作文也能用上，一举两得！

好了，准备好了吗？让我们一起穿越回春秋时代，去看看那些精彩的历史故事吧！

\newpage

\section*{\textcolor{orange!80!black}{一、郑伯克段于鄢}}

本篇讲述郑国国君庄公与他弟弟共叔段因为母亲姜氏偏爱而引发的争斗。庄公从小就因为难产不被母亲喜欢，而弟弟却备受宠爱。长大后，弟弟越来越贪婪，庄公却一直忍让。最后发生了什么呢？让我们一起来看这个经典的历史故事。

\vspace{0.3cm}
\sectab{\ruby{初}{chū}\refmark{1}，郑武公\ruby{娶}{qǔ}\refmark{2}于申，曰武姜\refmark{3}。生庄公及共叔段。庄公\ruby{寤}{wù}生\refmark{4}，惊姜氏，故名曰\ruby{寤}{wù}生，遂\ruby{恶}{wù}\refmark{5}之。爱共叔段，欲立之。\ruby{亟}{qì}\refmark{6}请于武公，公\ruby{弗}{fú}许\refmark{7}。}{最开始的时候，郑国的国君郑武公从申国娶了一位妻子，名叫武姜。她生了两个儿子：大儿子就是后来的庄公，小儿子叫共叔段。庄公出生的时候是脚先出来的，这是难产，把姜氏吓坏了。所以她给大儿子取名叫寤生，并且从此很讨厌这个儿子。她特别偏爱小儿子共叔段，想立他做太子，多次向武公请求，武公都没有答应。}{{\small \textcolor{blue}{(1)} 初：起初，当初。史书开头常用这个词来引出过去的某件事。}\\[0.1cm]{\small \textcolor{blue}{(2)} 娶：男子娶妻。}\\[0.1cm]{\small \textcolor{blue}{(3)} 武姜：武是她丈夫的谥号，姜是她娘家的姓。}\\[0.1cm]{\small \textcolor{blue}{(4)} 寤生：脚先出来的胎儿，就是难产。}\\[0.1cm]{\small \textcolor{blue}{(5)} 恶(wù)：厌恶，讨厌。}\\[0.1cm]{\small \textcolor{blue}{(6)} 亟(qì)：多次，屡次。}\\[0.1cm]{\small \textcolor{blue}{(7)} 弗许：不答应。}}

\vspace{0.3cm}
\sectab{及庄公即位，为之请制。公曰："制，岩邑也，\ruby{虢}{guó}叔\refmark{1}死\ruby{焉}{yān}\refmark{2}。佗邑唯命。"请京，使居之，谓之京城大叔。\ruby{祭}{zhài}仲\refmark{3}曰："都城过百\ruby{雉}{zhì}\refmark{4}，国之害也。先王之制：大都不过\ruby{参}{sān}\refmark{5}国之一；中，五之一；小，九之一。今京不度，非制也，君将不\ruby{堪}{kān}\refmark{6}。"公曰："姜氏欲之，\ruby{焉}{yān}\ruby{辟}{bì}\refmark{7}害？"对曰："姜氏何\ruby{厌}{yàn}\refmark{8}之有！不如早为之所，无使\ruby{滋}{zī}\ruby{蔓}{màn}\refmark{9}。\ruby{蔓}{màn}，难图也。蔓草犹不可除，况君之宠弟乎？"公曰："多行不义，必自\ruby{毙}{bì}\refmark{10}，子姑待之。"}{等到庄公当了国君之后，姜氏替共叔段向庄公请求封地，想要制这个地方。庄公说："制是地势险要的城邑，虢国的国君虢叔就死在那里。其他地方都可以听您的安排。"姜氏又请求京邑，庄公就让共叔段住在那里，人们都称他为京城大叔。郑国大夫祭仲对庄公说："都城的城墙如果超过一百雉，就会成为国家的祸害。按照先王的规定，大的都城不能超过国都的三分之一，中等的不能超过五分之一，小的不能超过九分之一。现在京邑的规模不合规矩，您将没办法控制。"庄公说："姜氏要这样，我能怎么办呢？"祭仲回答说："姜氏哪里会满足！不如早点安排处置他，不要让他发展壮大。一旦蔓延起来，就难对付了。蔓延的野草尚且除不掉，何况是您受宠的弟弟呢？"庄公说："多做不义的事情，一定会自取灭亡，您姑且等着看吧。"}{{\small \textcolor{blue}{(1)} 虢(guó)叔：虢国的国君。}\\[0.1cm]{\small \textcolor{blue}{(2)} 焉(yān)：在这里相当于"那里"，是一个代词。}\\[0.1cm]{\small \textcolor{blue}{(3)} 祭(zhài)仲：郑国大夫，一个很有眼光的大臣。}\\[0.1cm]{\small \textcolor{blue}{(4)} 雉(zhì)：古代城墙的长度单位，长三丈高一丈为一雉。}\\[0.1cm]{\small \textcolor{blue}{(5)} 参(sān)：同数字"三"。}\\[0.1cm]{\small \textcolor{blue}{(6)} 堪(kān)：承受，忍受。"不堪"就是受不了。}\\[0.1cm]{\small \textcolor{blue}{(7)} 辟(bì)：同"避"，躲避。}\\[0.1cm]{\small \textcolor{blue}{(8)} 厌(yàn)：满足。成语"贪得无厌"里的"厌"也是这个意思。}\\[0.1cm]{\small \textcolor{blue}{(9)} 滋蔓(zī màn)：滋生蔓延，越来越严重。}\\[0.1cm]{\small \textcolor{blue}{(10)} 毙(bì)：倒下，灭亡。"多行不义必自毙"后来成了成语。}}

\vspace{0.3cm}
\sectab{既而大叔命西\ruby{鄙}{bǐ}\refmark{1}北\ruby{鄙}{bǐ}\ruby{贰}{èr}\refmark{2}于己。公子吕曰："国不堪\ruby{贰}{èr}，君将若之何？欲与大叔，臣请事之；若弗与，则请除之，无生民心。"公曰："无\ruby{庸}{yōng}\refmark{3}，将自及。"大叔又收\ruby{贰}{èr}以为己邑，至于\ruby{廪}{lǐn}延\refmark{4}。子封曰："可矣，厚将得众。"公曰："不义不\ruby{昵}{nì}\refmark{5}，厚将\ruby{崩}{bēng}\refmark{6}。"}{不久之后，太叔命令西部和北部边境地区的城邑同时听从于他一个人。公子吕对庄公说："国家不能忍受这种两边都听命的情况，您打算怎么办？如果您想把君位让给太叔，我就请求去侍奉他；如果不让，就请除掉他，不要让百姓产生二心。"庄公说："不用，他会自己招来祸患的。"太叔又收取那些两属的地方作为自己的城邑，一直扩展到了廪延。子封说："可以动手了，他的势力大了就会得到更多人的拥护。"庄公说："对君不义，对兄不亲，势力再大也会崩溃。"}{{\small \textcolor{blue}{(1)} 鄙(bǐ)：边境上的城邑。}\\[0.1cm]{\small \textcolor{blue}{(2)} 贰(èr)：两属，从属于双方，既听庄公的也听太叔的。}\\[0.1cm]{\small \textcolor{blue}{(3)} 庸(yōng)：用。"无庸"就是"用不着"的意思。}\\[0.1cm]{\small \textcolor{blue}{(4)} 廪(lǐn)延：地名，在今河南延津附近。}\\[0.1cm]{\small \textcolor{blue}{(5)} 昵(nì)：亲近。"不昵"就是不肯亲近。}\\[0.1cm]{\small \textcolor{blue}{(6)} 崩(bēng)：崩溃，垮台。}}

\vspace{0.3cm}
\sectab{大叔完聚，\ruby{缮}{shàn}\refmark{1}\ruby{甲}{jiǎ}兵\refmark{2}，具卒\ruby{乘}{shèng}\refmark{3}，将\ruby{袭}{xí}\refmark{4}郑。夫人将启之。公闻其期，曰："可矣！"命子封\ruby{帅}{shuài}\refmark{5}车二百乘以\ruby{伐}{fá}京。京\ruby{叛}{pàn}\refmark{6}大叔段，段入于\ruby{鄢}{yān}\refmark{7}。公\ruby{伐}{fá}诸\ruby{鄢}{yān}。五月\ruby{辛}{xīn}\ruby{丑}{chǒu}，大叔出\ruby{奔}{bēn}\refmark{8}共。}{太叔加固城墙、聚集民众，修整铠甲和兵器，准备好了步兵和战车，准备偷袭郑国的都城。姜夫人也打算为他打开城门做内应。庄公得到了他们行动的日期，说："可以了！"于是命令子封率领二百辆战车去讨伐京邑。京邑的百姓背叛了太叔段，段逃到了鄢地。庄公又追到鄢地去讨伐他。五月辛丑这一天，太叔段逃出了郑国，投奔到共国去了。}{{\small \textcolor{blue}{(1)} 缮(shàn)：修整，修理。}\\[0.1cm]{\small \textcolor{blue}{(2)} 甲兵：铠甲和兵器。}\\[0.1cm]{\small \textcolor{blue}{(3)} 乘(shèng)：战车。古代四匹马拉一辆车为一乘。}\\[0.1cm]{\small \textcolor{blue}{(4)} 袭(xí)：乘人不备而进攻，就是偷袭。}\\[0.1cm]{\small \textcolor{blue}{(5)} 帅(shuài)：率领。同"率"字。}\\[0.1cm]{\small \textcolor{blue}{(6)} 叛(pàn)：背叛。}\\[0.1cm]{\small \textcolor{blue}{(7)} 鄢(yān)：地名，在今河南鄢陵。}\\[0.1cm]{\small \textcolor{blue}{(8)} 奔(bēn)：逃跑，逃亡。}}

\vspace{0.3cm}
\sectab{书曰："郑伯克段于\ruby{鄢}{yān}\refmark{1}。"段不弟\refmark{2}，故不言弟；如二君，故曰克\refmark{3}；称郑伯，讥失教\refmark{4}也；谓之郑志\refmark{5}。不言出奔，难之也。}{《春秋》上写道："郑伯克段于鄢。"因为段不遵守做弟弟的本分，所以不称他为"弟"；就像两个国君在打仗一样，所以用了"克"字；称庄公为"郑伯"，是批评他没有尽到教育弟弟的责任；说这是庄公的本意。不写"出奔"，是因为史官下笔的时候也很为难。}{{\small \textcolor{blue}{(1)} 鄢(yān)：地名。}\\[0.1cm]{\small \textcolor{blue}{(2)} 不弟：不遵守做弟弟的本分。}\\[0.1cm]{\small \textcolor{blue}{(3)} 克：战胜。用"克"字而不用其他字，当中含有微妙的批评意味，这就是"春秋笔法"。}\\[0.1cm]{\small \textcolor{blue}{(4)} 讥失教：批评郑伯没有教育好弟弟。}\\[0.1cm]{\small \textcolor{blue}{(5)} 郑志：郑伯的本意。}}

\vspace{0.3cm}
\sectab{遂\ruby{置}{zhì}\refmark{1}姜氏于城\ruby{颍}{yǐng}\refmark{2}，而\ruby{誓}{shì}\refmark{3}之曰："不及黄泉，无相见也。"既而悔之。\ruby{颍}{yǐng}考叔为\ruby{颍}{yǐng}谷封人，闻之，有献于公。公赐之食。食舍肉。公问之，对曰："小人有母，皆尝小人之食矣，未尝君之羹，请以\ruby{遗}{wèi}\refmark{4}之。"公曰："尔有母\ruby{遗}{wèi}，\ruby{繄}{yī}\refmark{5}我独无！"\ruby{颍}{yǐng}考叔曰："敢问何谓也？"公语之故，且告之悔。对曰："君何患\ruby{焉}{yān}？若\ruby{阙}{jué}\refmark{6}地及泉，\ruby{隧}{suì}\refmark{7}而相见，其谁曰不然？"公从之。}{于是庄公把姜氏安置在城颍这个地方，并且发誓说："不到黄泉，再也不见面！"但不久他就后悔了。颍考叔在颍谷做管理疆界的官员，听说了这件事，就找个机会向庄公进献礼物。庄公赐给他食物吃。他吃饭的时候把肉放在一边不吃。庄公问他为什么，他回答说："我有母亲，我的食物她都尝过了，但是没有尝过国君您的肉汤，请让我带回去给她。"庄公听了很伤感地说："你有母亲可以送东西给她，我偏偏没有！"颍考叔说："请问您说的是什么意思？"庄公告诉了他原因，也说出了自己的后悔。颍考叔回答说："您何必担心呢？如果挖地挖到泉水，在隧道里见面，谁能说这不是黄泉相见呢？"庄公听从了他的建议。}{{\small \textcolor{blue}{(1)} 置(zhì)：安置，安排。}\\[0.1cm]{\small \textcolor{blue}{(2)} 颍(yǐng)：地名，在今河南临颍西北。}\\[0.1cm]{\small \textcolor{blue}{(3)} 誓(shì)：发誓。}\\[0.1cm]{\small \textcolor{blue}{(4)} 遗(wèi)：赠送，给。}\\[0.1cm]{\small \textcolor{blue}{(5)} 繄(yī)：语气词，相当于"唯独"。}\\[0.1cm]{\small \textcolor{blue}{(6)} 阙(jué)：同"掘"，挖土。}\\[0.1cm]{\small \textcolor{blue}{(7)} 隧(suì)：隧道，地道。}}

\sectab{公入而\ruby{赋}{fù}\refmark{1}："大\ruby{隧}{suì}之中，其乐也\ruby{融}{róng}\ruby{融}{róng}\refmark{2}！"姜出而\ruby{赋}{fù}："大\ruby{隧}{suì}之外，其乐也\ruby{泄}{yì}\ruby{泄}{yì}\refmark{3}！"遂为母子如初。\par
君子曰："\ruby{颍}{yǐng}考叔，纯孝也，爱其母，\ruby{施}{yì}\refmark{4}及庄公。《诗》曰：'孝子不\ruby{匮}{kuì}\refmark{5}，永\ruby{锡}{cì}\refmark{6}尔类。'其是之谓乎？"}{庄公走进隧道，赋诗说："走进大隧道，心里真快乐！"姜氏走出隧道，也赋诗说："走出大隧道，心里真舒畅！"于是母子俩和好如初。君子评论说："颍考叔可算是真正的孝子了，他爱自己的母亲，并且把这份孝心扩大到了庄公身上。《诗经》上说：孝子的孝心没有穷尽，永远赐给你同类的人。说的就是这种情况吧？"}{{\small \textcolor{blue}{(1)} 赋(fù)：朗诵，吟咏诗句。}\\[0.1cm]{\small \textcolor{blue}{(2)} 融融：和睦快乐的样子。}\\[0.1cm]{\small \textcolor{blue}{(3)} 泄(yì)泄：舒畅快乐的样子。}\\[0.1cm]{\small \textcolor{blue}{(4)} 施(yì)及：扩展，影响到别人。}\\[0.1cm]{\small \textcolor{blue}{(5)} 匮(kuì)：穷尽，缺乏。}\\[0.1cm]{\small \textcolor{blue}{(6)} 锡(cì)：同"赐"，赐予。}}

\vspace{0.3cm}
{\large\textcolor{orange!80!black}{\textbf{想一想，说一说：}}}
\begin{enumerate}
  \item 庄公明明知道弟弟在扩张势力，为什么一直不制止他？
  \item 颍考叔用了一个什么巧妙的办法让庄公和母亲重新见面？
  \item "多行不义必自毙"这个成语是什么意思？你能用这个成语造个句子吗？
  \item 你觉得庄公是一个什么样的人？说说你的理由。
\end{enumerate}

\newpage

\section*{\textcolor{orange!80!black}{二、烛之武退秦师}}

本篇讲的是郑国被秦国和晋国两个大国围攻，眼看就要亡国了。在这危急关头，一个叫烛之武的老人，凭着一张能说会道的嘴，居然说服了秦国退兵。他是怎么做到的呢？让我们一起来看看古代外交家的智慧！

\vspace{0.3cm}
\sectab{晋\ruby{侯}{hóu}\refmark{1}、秦伯\ruby{围}{wéi}\refmark{2}郑，以其无礼于晋，且\ruby{贰}{èr}\refmark{3}于楚也。晋军\ruby{函}{hán}\ruby{陵}{líng}\refmark{4}，秦军\ruby{氾}{fán}\refmark{5}南。}{晋文公和秦穆公联合出兵围攻郑国。之所以打郑国，有两个原因：一是郑国曾经对晋文公无礼过，二是郑国对晋国怀有二心，暗中亲近楚国。晋国的军队驻扎在函陵，秦国的军队驻扎在氾水的南面。大军压境，郑国非常危险。}{{\small \textcolor{blue}{(1)} 侯(hóu)：诸侯的爵位，这里指晋文公。}\\[0.1cm]{\small \textcolor{blue}{(2)} 围(wéi)：围攻，包围。}\\[0.1cm]{\small \textcolor{blue}{(3)} 贰(èr)：对……有二心，暗中亲近。}\\[0.1cm]{\small \textcolor{blue}{(4)} 函(hán)陵：地名，在今河南新郑北面。}\\[0.1cm]{\small \textcolor{blue}{(5)} 氾(fán)：水名。}}

\vspace{0.3cm}
\sectab{\ruby{佚}{yì}之\ruby{狐}{hú}\refmark{1}言于郑伯曰："国危矣，若使烛之武\refmark{2}见秦君，师必退。"公从之。辞曰："臣之壮\refmark{3}也，犹不如人；今老矣，无能为也已。"公曰："吾不能早用子，今急而求子，是寡人\refmark{4}之过也。然郑亡，子亦有不利焉！"许之。}{郑国大夫佚之狐对郑文公说："国家危险了！如果派烛之武去面见秦君，秦军一定会撤退。"郑文公听从了他的建议。烛之武推辞说："我壮年的时候，尚且不如别人；现在老了，更是无能为力了。"郑文公说："我没能及早重用您，现在情况紧急了才来求您，这是我的过错。但是郑国灭亡了，对您也不利啊！"于是烛之武答应了。}{{\small \textcolor{blue}{(1)} 佚(yì)之狐：郑国大夫。}\\[0.1cm]{\small \textcolor{blue}{(2)} 烛之武：郑国大夫，这个故事的主角。}\\[0.1cm]{\small \textcolor{blue}{(3)} 壮：壮年，指年轻力壮的时候。}\\[0.1cm]{\small \textcolor{blue}{(4)} 寡人：古代国君对自己的谦称，意思是寡德之人。}}

\vspace{0.3cm}
\sectab{夜\ruby{缒}{zhuì}\refmark{1}而出。见秦伯曰："秦、晋围郑，郑既知亡矣。若亡郑而有益于君，敢以烦执事。越国以\ruby{鄙}{bǐ}\refmark{2}远，君知其难也。焉用亡郑以\ruby{陪}{péi}\refmark{3}邻？邻之厚，君之薄也。若舍郑以为东道主\refmark{4}，行李之往来，共其乏困，君亦无所害。且君尝为晋君赐矣，许君焦、\ruby{瑕}{xiá}\refmark{5}，朝\ruby{济}{jì}\refmark{6}而夕设\ruby{版}{bǎn}\refmark{7}焉，君之所知也。夫晋，何厌之有？既东封郑，又欲\ruby{肆}{sì}\refmark{8}其西封。若不\ruby{阙}{quē}\refmark{9}秦，将焉取之？阙秦以利晋，唯君图之。"}{当天夜里，烛之武用绳子从城墙上吊下来出了城。他见到秦穆公说："秦晋两国围攻郑国，郑国已经知道自己要灭亡了。如果灭亡郑国对您有好处，我怎敢拿这件事来麻烦您呢？可是，越过一个国家去把远方的地方作为边邑，您一定知道这是很困难的。既然如此，为什么要用灭亡郑国来增加您的邻国的实力呢？邻国实力增强了，就等于您的实力削弱了。如果您放过郑国，让它成为东方道路上的主人，贵国使者来往经过时，可以供应他们缺少的东西，这对您也没什么坏处。况且，您曾经给过晋国国君恩惠，他答应把焦、瑕两座城送给您，可他早上渡过黄河回国，晚上就筑城防御了，这是您知道的。晋国哪里会有满足的时候呢？已经在东面把郑国当作边界了，又想要扩张西面的边界。如果不损害秦国，它到哪里去夺取土地呢？损害秦国来让晋国得利，希望您好好考虑这件事。"}{{\small \textcolor{blue}{(1)} 缒(zhuì)：用绳子系着从城墙上吊下来。}\\[0.1cm]{\small \textcolor{blue}{(2)} 鄙(bǐ)：把……作为边境。}\\[0.1cm]{\small \textcolor{blue}{(3)} 陪(péi)：增加。}\\[0.1cm]{\small \textcolor{blue}{(4)} 东道主：东方道路上的主人。后来变成成语，指招待客人的主人。}\\[0.1cm]{\small \textcolor{blue}{(5)} 焦、瑕(xiá)：晋国的两座城邑。}\\[0.1cm]{\small \textcolor{blue}{(6)} 济(jì)：渡河。}\\[0.1cm]{\small \textcolor{blue}{(7)} 版(bǎn)：筑城墙用的夹板，这里指防御工事。}\\[0.1cm]{\small \textcolor{blue}{(8)} 肆(sì)：扩张，延伸。}\\[0.1cm]{\small \textcolor{blue}{(9)} 阙(quē)秦：损害秦国。}}

\vspace{0.3cm}
\sectab{秦伯说，与郑人\ruby{盟}{méng}\refmark{1}。使\ruby{杞}{qǐ}子\refmark{2}、逢孙、杨孙\ruby{戍}{shù}\refmark{3}之，乃还。}{秦穆公听了非常高兴，就和郑国签订了盟约。他还派杞子、逢孙、杨孙三位将领驻守在郑国帮助防守，自己就率领大军回国了。}{{\small \textcolor{blue}{(1)} 说(yuè)：同"悦"，高兴。}\\[0.1cm]{\small \textcolor{blue}{(2)} 盟(méng)：签订盟约。}\\[0.1cm]{\small \textcolor{blue}{(3)} 杞(qǐ)子：秦国大夫。}\\[0.1cm]{\small \textcolor{blue}{(4)} 戍(shù)：驻守，守卫。}}

\vspace{0.3cm}
\sectab{子\ruby{犯}{fàn}\refmark{1}请击之。公曰："不可。微\refmark{2}夫人\refmark{3}之力不及此。因人之力而\ruby{敝}{bì}\refmark{4}之，不仁；失其所与，不\ruby{知}{zhì}\refmark{5}；以乱易整，不\ruby{武}{wǔ}\refmark{6}。吾其还也。"亦去之。}{晋国大夫子犯请求追击秦军。晋文公说："不可以。没有那个人的帮助，我到不了今天这个地位。依靠了别人的力量却去损害他，这是不仁；失去自己的同盟者，这是不明智；用分裂来代替团结，这是不威武。我们还是回去吧。"于是晋军也撤兵了。烛之武凭着自己的智慧和口才，不费一兵一卒就解救了郑国的危机。}{{\small \textcolor{blue}{(1)} 犯(fàn)：狐偃，字子犯，是晋文公的舅舅。}\\[0.1cm]{\small \textcolor{blue}{(2)} 微：如果没有。}\\[0.1cm]{\small \textcolor{blue}{(3)} 夫人：那个人，指秦穆公。}\\[0.1cm]{\small \textcolor{blue}{(4)} 敝(bì)：损害，伤害。}\\[0.1cm]{\small \textcolor{blue}{(5)} 知(zhì)：同"智"，明智。}\\[0.1cm]{\small \textcolor{blue}{(6)} 武(wǔ)：威武，合乎武德。}}

\vspace{0.3cm}
{\large\textcolor{orange!80!black}{\textbf{想一想，说一说：}}}
\begin{enumerate}
  \item 烛之武一开始不愿意去见秦伯，这是为什么？
  \item 烛之武用了哪些理由来说服秦伯退兵的？请列举出来。
  \item "东道主"这个成语是怎么来的？现在在生活中是什么意思？
  \item 你觉得烛之武是一个什么样的人？你最佩服他哪一点？
\end{enumerate}

\newpage

\section*{\textcolor{orange!80!black}{三、曹刿论战}}

这篇故事讲了鲁国和齐国之间的一场战争。鲁国比齐国弱小，但却打赢了。为什么呢？因为一个叫曹刿的谋士帮鲁庄公出谋划策。他提出了一个著名的理论——"一鼓作气"。这个成语我们现在还在用呢！

\vspace{0.3cm}
\sectab{十年春，齐师\ruby{伐}{fá}\refmark{1}我。公将战，曹\ruby{刿}{guì}\refmark{2}请见。其乡人\refmark{3}曰："肉食者\refmark{4}\ruby{谋}{móu}之，又何\ruby{间}{jiàn}\refmark{5}焉？"\ruby{刿}{guì}曰："肉食者\ruby{鄙}{bǐ}\refmark{6}，未能远\ruby{谋}{móu}。"乃入见。问："何以战？"公曰："衣食所\ruby{安}{ān}，弗敢专也，必以分人。"对曰："小\ruby{惠}{huì}未\ruby{遍}{biàn}\refmark{7}，民弗从也。"公曰："牺牲玉帛\refmark{8}，弗敢加也，必以信。"对曰："小信未\ruby{孚}{fú}\refmark{9}，神弗福也。"公曰："小大之\ruby{狱}{yù}\refmark{10}，虽不能察，必以情。"对曰："忠之属也，可以一战。战则请从。"}{鲁庄公十年的春天，齐国的军队攻打鲁国。鲁庄公准备迎战。曹刿请求拜见庄公。他的同乡说："当权的人会谋划这件事的，你又何必参与呢？"曹刿说："当权的人目光短浅，不能深谋远虑。"于是入宫进见。他问庄公："您凭借什么来作战？"庄公说："衣食这类用来安身的东西，我不敢独自享受，一定分给别人。"曹刿说："这些小恩惠不能遍及百姓，人民是不会跟从您的。"庄公说："祭祀用的牛羊玉帛之类，我不敢虚报，一定诚实。"曹刿说："这只是小信用，不能使神灵信服。"庄公说："大大小小的案件，我虽然不能一一查明，但一定根据实情来判断。"曹刿说："这才是尽了君主的本分，可以凭借这一点打一仗。作战时请让我跟您一起去。"}{{\small \textcolor{blue}{(1)} 伐(fá)：攻打。}\\[0.1cm]{\small \textcolor{blue}{(2)} 刿(guì)：曹刿，鲁国的谋士。}\\[0.1cm]{\small \textcolor{blue}{(3)} 乡人：同乡的人。}\\[0.1cm]{\small \textcolor{blue}{(4)} 肉食者：吃肉的人，指当权的官员。}\\[0.1cm]{\small \textcolor{blue}{(5)} 间(jiàn)：参与。}\\[0.1cm]{\small \textcolor{blue}{(6)} 鄙(bǐ)：目光短浅。}\\[0.1cm]{\small \textcolor{blue}{(7)} 遍(biàn)：遍及，普遍。}\\[0.1cm]{\small \textcolor{blue}{(8)} 牺牲玉帛：祭祀用的祭品。}\\[0.1cm]{\small \textcolor{blue}{(9)} 孚(fú)：使人信服。}\\[0.1cm]{\small \textcolor{blue}{(10)} 狱(yù)：案件。}}

\vspace{0.3cm}
\sectab{公与之\ruby{乘}{chéng}\refmark{1}，战于长\ruby{勺}{sháo}\refmark{2}。公将\ruby{鼓}{gǔ}\refmark{3}之。\ruby{刿}{guì}曰："未可。"齐人三鼓。\ruby{刿}{guì}曰："可矣。"齐师败\ruby{绩}{jì}\refmark{4}。公将\ruby{驰}{chí}\refmark{5}之。\ruby{刿}{guì}曰："未可。"下视其\ruby{辙}{zhé}\refmark{6}，登\ruby{轼}{shì}\refmark{7}而望之，曰："可矣。"遂逐齐师。}{鲁庄公和曹刿同乘一辆战车，在长勺这个地方与齐军交战。庄公刚要击鼓进军，曹刿说："还不行。"等到齐军击了三次鼓之后，曹刿说："可以了。"齐军大败。庄公要驱车追击，曹刿又说："还不行。"他下车仔细看了齐军战车留下的痕迹，又登上车前横木远望齐军队形，然后说："可以了。"于是下令追击齐军，取得了胜利。}{{\small \textcolor{blue}{(1)} 乘(chéng)：乘坐战车。}\\[0.1cm]{\small \textcolor{blue}{(2)} 长勺(sháo)：鲁国地名，在今山东莱芜东北。}\\[0.1cm]{\small \textcolor{blue}{(3)} 鼓(gǔ)：击鼓，古代击鼓是进攻的信号。}\\[0.1cm]{\small \textcolor{blue}{(4)} 败绩(jì)：大败。}\\[0.1cm]{\small \textcolor{blue}{(5)} 驰(chí)：驱车追赶。}\\[0.1cm]{\small \textcolor{blue}{(6)} 辙(zhé)：车轮碾过的痕迹。}\\[0.1cm]{\small \textcolor{blue}{(7)} 轼(shì)：古代车厢前用作扶手的横木。}}

\vspace{0.3cm}
\sectab{既克，公问其故。对曰："夫战，勇气也。一鼓作气，再而\ruby{衰}{shuāi}\refmark{1}，三而\ruby{竭}{jié}\refmark{2}。彼竭我\ruby{盈}{yíng}\refmark{3}，故克之。夫大国，难\ruby{测}{cè}\refmark{4}也，\ruby{惧}{jù}有\ruby{伏}{fú}\refmark{5}焉。吾视其辙乱，望其旗\ruby{靡}{mǐ}\refmark{6}，故逐之。"}{战胜之后，庄公问曹刿这样做的原因。曹刿回答说："作战，靠的是勇气。第一次击鼓能够振作士兵的勇气，第二次击鼓勇气就衰弱了，第三次击鼓勇气就消耗完了。对方的勇气耗尽了，而我方的勇气正饱满，所以战胜了他们。大国是难以预测的，我担心他们有埋伏。后来我看到他们的车辙混乱，远望他们的旗帜也倒下了，知道他们是真败，所以下令追击。"}{{\small \textcolor{blue}{(1)} 衰(shuāi)：衰弱，减弱。}\\[0.1cm]{\small \textcolor{blue}{(2)} 竭(jié)：枯竭，耗尽。}\\[0.1cm]{\small \textcolor{blue}{(3)} 盈(yíng)：充满，旺盛。}\\[0.1cm]{\small \textcolor{blue}{(4)} 测(cè)：预测，推测。}\\[0.1cm]{\small \textcolor{blue}{(5)} 伏(fú)：埋伏。}\\[0.1cm]{\small \textcolor{blue}{(6)} 靡(mǐ)：倒下。}}

\vspace{0.3cm}
{\large\textcolor{orange!80!black}{\textbf{想一想，说一说：}}}
\begin{enumerate}
  \item 曹刿为什么认为鲁庄公可以打这一仗？
  \item "一鼓作气"是什么意思？曹刿为什么要等齐人三鼓之后才出击？
  \item 曹刿在追击之前为什么要先看齐军的车辙和旗帜？
  \item 从这个故事中你学到了什么道理？
\end{enumerate}

\newpage

\section*{\textcolor{orange!80!black}{四、弦高犒师}}

这篇故事的主人公不是什么大英雄大人物，而是一个普普通通的商人。他在去外地做生意的路上，偶然发现了一个巨大的阴谋——有人要偷袭他的国家！他急中生智，用自己的货物巧妙地化解了一场战争。一起去看看这位爱国商人的故事吧！

\vspace{0.3cm}
\sectab{三十三年春，秦师过周北门，左右免\ruby{胄}{zhòu}\refmark{1}而下，超乘\refmark{2}者三百乘。王孙满尚幼，观之，言于王曰："秦师轻\refmark{3}而无礼，必败。轻则寡\ruby{谋}{móu}，无礼则\ruby{脱}{tuō}\refmark{4}。入险而脱，又不能谋，能无败乎？"}{鲁僖公三十三年春天，秦国的军队经过周朝都城的北门。战车上的左右卫士虽然脱下了头盔下车，但有三百辆战车的士兵又轻狂地跳上车去。当时周王室的子孙满年纪还小，看到这个情景，对周王说："秦国的军队轻佻而无礼，一定会失败。轻佻就会缺少谋略，无礼就会粗心大意。进入险境还粗心大意，又没有谋略，能不失败吗？"}{{\small \textcolor{blue}{(1)} 胄(zhòu)：头盔。}\\[0.1cm]{\small \textcolor{blue}{(2)} 超乘：跳跃着上车。这种举动显得轻狂无礼。}\\[0.1cm]{\small \textcolor{blue}{(3)} 轻：轻佻，不庄重。}\\[0.1cm]{\small \textcolor{blue}{(4)} 脱(tuō)：粗心，疏忽。}}

\vspace{0.3cm}
\sectab{及滑，郑商人弦高\refmark{1}将市于周，遇之。以乘韦\refmark{2}先，牛十二\ruby{犒}{kào}\refmark{3}师，曰："寡君闻吾子将步师出于敝邑，敢\ruby{犒}{kào}从者。不\ruby{腆}{tiǎn}\refmark{4}敝邑，为从者之淹，居则具一日之积，行则备一夕之卫。"且使\ruby{遽}{jù}\refmark{5}告于郑。}{秦军到了滑国的时候，郑国商人弦高正要去周地做买卖，正好遇上了秦军。弦高先送四张熟牛皮，再送十二头牛去犒劳秦军，说："我们国君听说您要行军到我们国家去，特地派我来犒劳您的部下。我们国家虽然不富裕，但为了您的部下在这里停留，住下就准备好一天的粮草，离开就准备好一夜的守卫。"弦高说这话的同时，已经派人骑着快马紧急向郑国报告去了。}{{\small \textcolor{blue}{(1)} 弦高：郑国的商人，一位爱国的义士。}\\[0.1cm]{\small \textcolor{blue}{(2)} 乘韦：四张熟牛皮。"乘"代指"四"，"韦"指熟牛皮。}\\[0.1cm]{\small \textcolor{blue}{(3)} 犒(kào)：用酒食等慰劳。}\\[0.1cm]{\small \textcolor{blue}{(4)} 腆(tiǎn)：丰厚。"不腆"就是不丰厚。}\\[0.1cm]{\small \textcolor{blue}{(5)} 遽(jù)：急忙，迅速。}}

\vspace{0.3cm}
\sectab{郑穆公使视客馆，则束载、厉兵、\ruby{秣}{mò}\refmark{1}马矣。使皇武子辞焉，曰："吾子淹久于敝邑，唯是\ruby{脯}{fǔ}资\ruby{饩}{xì}牵\refmark{2}竭矣。为吾子之将行也，郑之有原圃\refmark{3}，犹秦之有具囿\refmark{4}也，吾子取其麋鹿，以闲敝邑，若何？"\ruby{杞}{qǐ}子\refmark{5}奔齐，逢孙、杨孙奔宋。}{郑穆公接到弦高的报告，派人到秦国驻军的客馆去查看，发现他们已经捆好行装、磨好了兵器、喂饱了战马。于是郑穆公派皇武子去下逐客令，说："你们在这里住得太久了，粮食物资都已经用完了。郑国有打猎的原圃，就像秦国有打猎的具囿一样，你们自己去打些麋鹿来补充吧。"秦国的将领们知道阴谋已经败露，杞子逃往齐国，逢孙和杨孙逃往宋国。}{{\small \textcolor{blue}{(1)} 秣(mò)马：喂饱战马。成语"厉兵秣马"就出自这里。}\\[0.1cm]{\small \textcolor{blue}{(2)} 脯(fǔ)资饩(xì)牵：干肉、粮食和活牲畜等物资。}\\[0.1cm]{\small \textcolor{blue}{(3)} 原圃：郑国养禽兽的猎苑。}\\[0.1cm]{\small \textcolor{blue}{(4)} 具囿(yòu)：秦国养禽兽的猎苑。}\\[0.1cm]{\small \textcolor{blue}{(5)} 杞(qǐ)子：秦国的大夫。}}

\vspace{0.3cm}
\sectab{孟明\refmark{1}曰："郑有备矣，不可\ruby{冀}{jì}\refmark{2}也。攻之不克，围之不继，吾其还也。"灭滑而还。}{秦国的主帅孟明说："郑国已经有防备了，不能指望偷袭成功了。要攻打攻不下来，要包围又没有后续部队，我们还是回去吧。"于是秦军顺便灭了滑国就回去了。就这样，一个普通的商人弦高，凭着自己的机智和爱国精神，避免了一场战争，保护了自己的国家。}{{\small \textcolor{blue}{(1)} 孟明：秦国的大将。}\\[0.1cm]{\small \textcolor{blue}{(2)} 冀(jì)：希望，指望。}}

\vspace{0.3cm}
{\large\textcolor{orange!80!black}{\textbf{想一想，说一说：}}}
\begin{enumerate}
  \item 弦高只是一个普通的商人，他为什么还要管国家的事？
  \item 弦高用了什么办法让秦军以为郑国已经有防备了？
  \item 如果弦高没有出手相助，郑国可能会怎样？
  \item 你觉得"天下兴亡，匹夫有责"这句话和这个故事有关系吗？
\end{enumerate}

\newpage

\section*{\textcolor{orange!80!black}{五、唇亡齿寒}}

这个故事讲了一个非常有名的成语——"唇亡齿寒"是怎么来的。晋国想攻打虢国，向虞国借路。虞国的大子宫之奇看出了晋国的阴谋，苦苦劝说虞公不要借路。可是虞公就是不听。最后发生了什么？这个成语又是什么意思呢？

\vspace{0.3cm}
\sectab{晋侯复假道于\ruby{虞}{yú}\refmark{1}以\ruby{伐}{fá}\ruby{虢}{guó}\refmark{2}。宫之奇\refmark{3}谏曰："\ruby{虢}{guó}，虞之表\refmark{4}也。\ruby{虢}{guó}亡，虞必从之。晋不可启，\ruby{寇}{kòu}\refmark{5}不可\ruby{翫}{wán}\refmark{6}。一之谓甚，其可再乎？\ruby{谚}{yàn}\refmark{7}所谓'辅车相依，唇亡齿寒'者，其虞、\ruby{虢}{guó}之谓也。"}{晋献公再次向虞国借路，要去攻打虢国。虞国大夫宫之奇劝阻虞公说："虢国是虞国的外围屏障。虢国如果灭亡了，虞国也一定跟着灭亡。晋国的野心不能随便开启，国外的敌人不能疏忽大意。借一次路已经够过分了，怎么还能再借第二次呢？俗话说辅车相依、唇亡齿寒，说的就是虞国和虢国这种互相依存的关系啊！"}{{\small \textcolor{blue}{(1)} 虞(yú)：诸侯国名。}\\[0.1cm]{\small \textcolor{blue}{(2)} 虢(guó)：诸侯国名。}\\[0.1cm]{\small \textcolor{blue}{(3)} 宫之奇：虞国的大夫，一位有远见的政治家。}\\[0.1cm]{\small \textcolor{blue}{(4)} 表：屏障，外围。}\\[0.1cm]{\small \textcolor{blue}{(5)} 寇(kòu)：外敌，侵略者。}\\[0.1cm]{\small \textcolor{blue}{(6)} 翫(wán)：同"玩"，忽视，轻慢。}\\[0.1cm]{\small \textcolor{blue}{(7)} 谚(yàn)：谚语，俗语。}}

\vspace{0.3cm}
\sectab{公曰："晋，吾宗也，岂害我哉？"对曰："大伯、虞仲，大王之\ruby{昭}{zhāo}也。大伯不从，是以不\ruby{嗣}{sì}\refmark{1}。\ruby{虢}{guó}仲、\ruby{虢}{guó}叔，王季之穆也，为文王卿士，\ruby{勋}{xūn}在王室，藏于盟府。将\ruby{虢}{guó}是灭，何爱于虞？且虞能亲于桓、庄乎，其爱之也？桓、庄之族何罪，而以为\ruby{戮}{lù}\refmark{2}，不唯逼\refmark{3}乎？亲以宠逼，犹尚害之，况以国乎？"}{虞公说："晋国和我是同宗，难道会害我吗？"宫之奇回答说："虢国也是晋国的同宗，晋国连虢国都要灭掉，对虞国又有什么爱惜呢？况且，晋国公室中桓叔和庄伯的家族有什么罪过？全被杀害了，不就是因为他们势力太大威胁到了国君吗？至亲的人因为受宠而威胁到国君，尚且被杀害，何况您有一个国家这么大的威胁呢？"}{{\small \textcolor{blue}{(1)} 嗣(sì)：继承。}\\[0.1cm]{\small \textcolor{blue}{(2)} 戮(lù)：杀戮。}\\[0.1cm]{\small \textcolor{blue}{(3)} 逼：同"逼"，威胁。}}

\vspace{0.3cm}
\sectab{公曰："吾享祀丰洁，神必据我。"对曰："臣闻之，鬼神非人实亲，惟德是依。故《周书》曰：'皇天无亲，惟德是辅。'又曰：'\ruby{黍}{shǔ}\ruby{稷}{jì}\refmark{1}非\ruby{馨}{xīn}\refmark{2}，明德惟\ruby{馨}{xīn}。'如是，则非德，民不和，神不享矣。神所冯依，将在德矣。"}{虞公又说："我祭祀的祭品丰盛洁净，神灵一定会保佑我的。"宫之奇回答说："我听说，鬼神并不亲近任何人，只帮助有德行的人。所以《周书》上说：上天没有亲疏之分，只帮助有德行的人。又说：黍稷并不是芳香，美德才是芳香。照这样说，如果没有德行，人民就不会和睦，神灵也不会享用祭品。神灵所依靠的，就在于德行了。"}{{\small \textcolor{blue}{(1)} 黍稷(shǔ jì)：黍和稷，泛指五谷，用作祭祀的供品。}\\[0.1cm]{\small \textcolor{blue}{(2)} 馨(xīn)：芳香，这里指祭品的香气。}}

\sectab{弗听，许晋使。宫之奇以其族行，曰："虞不\ruby{腊}{là}\refmark{1}矣，在此行也，晋不更举矣。"\par
冬，十二月丙子\ruby{朔}{shuò}\refmark{2}，晋灭\ruby{虢}{guó}。\ruby{虢}{guó}公丑奔京师。师还，馆于虞，遂袭虞，灭之。执虞公。}{虞公不听宫之奇的劝告，答应了晋国使者借路的请求。宫之奇带着全族的人离开了虞国，临走时说："虞国等不到年终祭祀的那一天了，就在这次行动中，晋国不需要再次出兵了。"果然，冬天十二月初一，晋国灭掉了虢国。晋军回国的时候，驻扎在虞国，趁机偷袭了虞国，灭了它，抓走了虞公。宫之奇的预言全都应验了。}{{\small \textcolor{blue}{(1)} 腊(là)：年终举行的一种祭祀。}\\[0.1cm]{\small \textcolor{blue}{(2)} 朔(shuò)：农历每月的第一天，也就是初一。}}

\vspace{0.3cm}
{\large\textcolor{orange!80!black}{\textbf{想一想，说一说：}}}
\begin{enumerate}
  \item "唇亡齿寒"这个成语的字面意思是什么？它比喻什么？
  \item 宫之奇用了哪些理由劝虞公不要借路？
  \item 虞公为什么不听宫之奇的劝告？
  \item 这个故事的结局说明了什么道理？
\end{enumerate}

\newpage

\section*{\textcolor{orange!80!black}{六、退避三舍}}

这篇故事讲的是晋文公重耳在流亡的时候，曾经答应楚王：如果将来两国交战，我会主动后退九十里来报答您。后来他当上了国君，两国真的打起来了。他会不会履行当初的诺言呢？"退避三舍"这个成语就是这么来的。

\vspace{0.3cm}
\sectab{晋公子重耳\refmark{1}之及于难，晋人\ruby{伐}{fá}诸蒲城。蒲城人欲战，重耳不可，曰："保君父之命而享其生禄，于是乎得人。有人而\ruby{校}{jiào}\refmark{2}，罪莫大焉。吾其奔也。"遂奔狄。}{晋国的公子重耳遭遇了迫害，原因是晋献公听信了骊姬的谗言要杀他。晋国的军队到蒲城去攻打他。蒲城的人想要迎战，重耳不同意，说："依靠君父的命令才能享有俸禄，才能得到百姓的拥护。有了百姓的拥护却去反抗君父，没有比这更大的罪过了。我还是逃亡吧。"于是逃到了狄国。}{{\small \textcolor{blue}{(1)} 重耳：晋献公的儿子，后来成为春秋五霸之一的晋文公。}\\[0.1cm]{\small \textcolor{blue}{(2)} 校(jiào)：同"较"，较量，反抗。}}

\vspace{0.3cm}
\sectab{及楚，楚子\ruby{飨}{xiǎng}\refmark{1}之，曰："公子若反晋国，则何以报不\ruby{穀}{gǔ}\refmark{2}？"对曰："子女玉帛，则君有之；羽毛齿革，则君地生焉。其波及晋国者，君之余也，其何以报？"曰："虽然，何以报我？"对曰："若以君之灵，得反晋国，晋、楚治兵，遇于中原，其\ruby{辟}{bì}\refmark{3}君三\ruby{舍}{shè}\refmark{4}。若不获命，其左执鞭\ruby{弭}{mǐ}\refmark{5}，右属\ruby{櫜}{gāo}\ruby{鞬}{jiān}\refmark{6}，以与君周旋。"}{重耳流亡到了楚国，楚成王设宴款待他，问道："公子如果回到晋国当上了国君，用什么来报答我呢？"重耳回答说："美女、宝玉和丝绸，您有的是；鸟羽、象牙和皮革，您的土地上也出产。那些流散到晋国的，不过是您剩下的罢了，我拿什么来报答您呢？"楚王说："话虽这么说，可你总得报答我吧？"重耳回答说："如果托您的福，我能回到晋国，将来晋楚两国交战，在中原相遇，我一定后退三舍来报答您。如果这样还得不到您的谅解，那么我就只好左手拿着鞭子和弓，右边挂着弓箭袋，跟您较量一番了。"}{{\small \textcolor{blue}{(1)} 飨(xiǎng)：用酒食款待，设宴招待。}\\[0.1cm]{\small \textcolor{blue}{(2)} 不穀(gǔ)：古代诸侯对自己的谦称，相当于"我"。}\\[0.1cm]{\small \textcolor{blue}{(3)} 辟(bì)：同"避"，退避，避开。}\\[0.1cm]{\small \textcolor{blue}{(4)} 舍(shè)：古代行军一舍为三十里，三舍就是九十里。}\\[0.1cm]{\small \textcolor{blue}{(5)} 弭(mǐ)：一种不加装饰的弓。}\\[0.1cm]{\small \textcolor{blue}{(6)} 櫜鞬(gāo jiān)：装弓箭的袋子。}}

\sectab{后三年，晋侯\refmark{1}围楚于城\ruby{濮}{pú}\refmark{2}。楚子使子玉去宋，曰："无从晋师。"子玉曰："王遇之至厚，今知子玉之未能。"\par
晋师退三\ruby{舍}{shè}\refmark{3}。楚众欲止，子玉不可。}{三年之后，晋文公在城濮包围了楚军。楚成王命令大将子玉不要去追击晋军。子玉不听。晋文公履行了当初的诺言，命令晋军后退九十里。楚军想要停止追击，但子玉不同意。最后，晋军在城濮之战中大败楚军，晋文公因此成为了春秋五霸之一。"退避三舍"这个成语也一直流传到了今天，用来比喻退让和回避。}{{\small \textcolor{blue}{(1)} 晋侯：晋文公重耳。}\\[0.1cm]{\small \textcolor{blue}{(2)} 濮(pú)：城濮，地名。著名的城濮之战就发生在这里。}\\[0.1cm]{\small \textcolor{blue}{(3)} 舍(shè)：三舍即九十里。成语"退避三舍"就出自这里，现在用来比喻主动退让，避免冲突。}}

\vspace{0.3cm}
{\large\textcolor{orange!80!black}{\textbf{想一想，说一说：}}}
\begin{enumerate}
  \item 重耳为什么要对楚王承诺退避三舍？
  \item "退避三舍"这个成语现在是什么意思？你能用它造个句子吗？
  \item 晋文公后来真的退避三舍了，这说明了什么品格？
\end{enumerate}

\vspace{0.5cm}
\begin{center}
\textcolor{blue!70!black}{\large —— 加油！你们都是最棒的！ ——}
\end{center}
\end{document}
